% Source-derived chapter 03: Sheaves on Hida families
% Original source: universal-p-adic-gross-zagier-formula
\chapter{Sheaves on Hida families}
\label{universal-p-adic-gross-zagier-formula-chapter-03}
\source{universal-p-adic-gross-zagier-formula}

\begin{remark}[Source section]
\label{universal-p-adic-gross-zagier-formula-chapter-03-source}
\source{universal-p-adic-gross-zagier-formula}
\source{universal-p-adic-gross-zagier-formula}
This chapter reproduces the corresponding section of the retrieved
LaTeX source, including its definitions, statements, examples, and
proofs. It has not yet been translated into Lean.
\notready
\end{remark}

% The source section title was: Sheaves on Hida families
\label{section 3}
\source{universal-p-adic-gross-zagier-formula}
We construct the universal Hecke- and Galois- modules over Hida families for $(\G\ts\H)'$ and prove a local-global compatibility result. We claim no originality for the  results of \S\S~\ref{3.1}-\ref{3.2new}.
\subsection{Hida theory}\lb{3.1}
We let   $\G_{*}$ denote any of the groups $\G$, $\H$, $(\G\times \H)$, $(\G\times \H)'$, and let $r\in \N^{\{v\vert p\}}$. We will use the notation from \S~\ref{sec: not}.
 For $U_{*}^{p}\subset \G(\A^{p\infty})$ we let $X_{*, U_{*}^{p},r}:=X_{*, U_{*}^{p}U_{*,r}}$ be the corresponding Shimura variety.









  
  When $M$ is a $\Z_{p}$-module with   action by $T_{\G_{*}}^{+}$, arising as limit of ordinary parts of   $p$-adic coadmissible  $\G_{*}(\Q_{p})\ts G_{*, \infty}$-modules (see Definition \ref{def coadmiss} and \S~\ref{ssec ord}),
  we denote this action by 
  $$t\mapsto \Up_{t}$$
  and adopt the notation of \eqref{aut-Up}. 







\subsubsection{Weight spaces} \lb{sec wtsp} Let $U_{*}^{p}\subset \G_{*}(\A^{p\infty})$ be a compact open subgroup, and define $Z_{\G_{*}, U^{p}_{*}}
\subset  {\rm Z}_{\G}(\Q)$ by 
\beqq
Z_{\G, U^{p}}&:= {\rm Z}_{\G}(\Q)\cap U^{p}T_{\G, 0}, 
&Z_{\H, V^{p}}&:=\H(\Q)\cap V^{p}T_{\H, 0}=\OO_{E}^{\times}\cap V^{ p }, & {\ } \\
Z_{(\G\times \H)', K^{p}}&:= \text{the image in $T_{(\G\times \H)'}$ of }  &Z_{\G\times H, U^{p}\times V^{p}} : =  Z_{\G, U^{p}}\times Z_{\H, V^{p}} & \text{if $K^{p}$ is the image of $U^{p}\times V^{p}$}.
\eeqq 
In all cases, let 
$\baar{T}_{\G_{*},U_{*}^{p},0}:=T_{\G,0}/\baar{Z_{\G_{*}, U_{*}}}$,
where $\baar{\Box}$  denotes the closure for the $p$-adic topology, and let  $ \baar{T}_{\G_{*}, U^{p}_{*}, r}\subset \baar{T}_{\G, U^{p}_{*},0}$ be the image of $T_{\G_{*},U^{p}_{*},r}$.
Let
$$\Lm^{\circ}_{\G_{*}, U^{p}_{*}}:=\Z_{p}\llb \baar{T}_{\G_{*},0}\rrb, 
$$
and for an irreducible algebraic representation $W$ of $\G_{*}$ consider the ideals 
\beq 
\label{Ir}
\source{universal-p-adic-gross-zagier-formula}
I_{\G_{*},U^{p}_{*}, W, r, L}:= ([t]-\sigma_{W}^{-1}(t))_{t\in \baar{T}_{\G_{*},  U^{p}_{*},r}})\subset \Lm^{\circ}_{\G_{*}  ,U^{p}_{*}}\otimes \OO_{L}
\eeq
For each fixed $W$ and varying $r$, the ideals $I_{\G_{*},U^{p}_{*}, W, r}$ form a fundamental system of neighbourhoods of zero in $\Lm^{\circ}_{\G_{*}, U^{p}_{*}}\otimes \OO_{L}$, so that  
\beq\label{alg density}
\source{universal-p-adic-gross-zagier-formula}
\Lm_{\G_{*}, U^{p}_{*}}:= \Lm^{\circ}_{\G_{*}, U^{p}_{*}}\otimes_{\OO_{L}} L=(\varprojlim_{r}\Lambda_{\G_{*}, U^{p}_{*}, W,r})\ot_{\OO_{L}} L\eeq
with
\beq\label{lambda r}
\source{universal-p-adic-gross-zagier-formula}
\Lambda^{\circ}_{\G_{*}, U^{p}_{*}, W, r}:= \Lambda^{\circ}_{\G_{*}, U^{p}_{*}}
/ I_{\G_{*}, U^{p}_{*},W, r}\cong
 \OO_{L}[\baar{T}_{\G_{*}, U^{p}_{*}, 0}/\baar{T}_{\G_{*}, U^{p}_{*}, r}],
\eeq
where the isomorphism is given by $[t]\mapsto \sigma_{W}^{-1}(t)[\baar{t}]$.
When $W=\Q_{p}$, we omit $W$ from the  notation. We also omit the subscript $U^{p}_{*}$  when it is unimportant or understood from context.


Writing $\baar{T}_{\G_{*},U^{p}_{*},0}\cong \Delta\times \Z_{p}^{{\rm d}(\G_{*})}$ for a finite torsion group $\Delta$, we have an isomorphism $\Lambda_{\G_{*}, U^{p}_{*}}\cong \Z_{p}[\Delta]\otimes  \Z_{p}\llb X_{1}, \ldots X_{{\rm d}(\G_{*})}\rrb$ for  an integer ${{\rm d}(\G_{*})}$  given by\footnote{We would have ${\rm d}((\G\ts_{\rm Z}\H)'=2[F:\Q]$ for the group of the footnote after \eqref{list}, which explains the way the number of variables of Theorem \ref{iw} is counted in the abstract.}
$${\rm d}(\G)= {\rm d}(\H)= [F:\Q]+1+\delta, \quad {\rm d}((\G\times \H)')= 2[F:\Q] + 1+\delta,$$
where $\delta=\delta_{F, p}$ is the Leopoldt defect of $F$ at $p$; see \cite[\S~2.2.3.3]{fouquet} for ${\rm d}(\G)$.

\begin{defi}
\source{universal-p-adic-gross-zagier-formula}
\notready \label{clwt} 
\source{universal-p-adic-gross-zagier-formula}
The \emph{weight space} is 
$$  \mathfrak{W}_{\G_{*}, U^{p}_{*}}:= \Spec \Lambda_{\G_{*}, U^{p}_{*}}\otimes \Q_{p}.$$
Let $W$ be an irreducible cohomological  algebraic representation of $\G_{*}$. The zero-dimensional subscheme of \emph{classical points of weight $W$}  and level $r$  is 
$$\mathfrak{W}_{\G_{*}, U^{p}_{*},r}^{\cl, W}:=  \Spec \Lambda_{\G_{*}, U^{p}_{*}, r, W}.$$
The ind-subschemes   of all classical points of weight $W$ and of of all classical points  are respectively
$$\mathfrak{W}_{\G_{*}, U^{p}_{*}}^{\cl, W}:=\bigcup_{r\geq 0}  \mathfrak{W}_{\G_{*}, U^{p}_{*},r}^{\cl, W}, \qquad\mathfrak{W}_{\G_{*}, U^{p}_{*}}^{\cl}:=\bigcup_{W} \mathfrak{W}_{\G_{*}, U^{p}_{*}}^{\cl, W},$$
where as usual the union runs through the algebraic representations of cohomological weight.
\end{defi}







\subsubsection{Ordinary completed homology} 
Let $W$ be an irreducible right   algebraic representation of $\G_{*}(\Q_{p})$ over $L$, and fix a $\G(\Z_{p})$-stable $\OO_{L}$-lattice $W^{\circ}\subset W$. 
Let $\cW$ be the  local system attached to $W$,  and for $U_{*}^{p}\subset \G_{*}(\A^{p\infty})$, $r\geq 0$ consider the \emph{ordinary parts}
\beqq\H^{\textrm{\'et}}_{d}(\baar{X}_{*,U_{*}^{p}U_{*,r}}, \cW^{\circ})^{\ord} &
 := (H^{\textrm{\'et}}_{d}(\baar{X}_{*,U_{*}^{p}U_{*,r}}, \cW^{\circ})\ot_{\OO_{L}[T^{+}_{\G_{*}}]} (W^{\circ,\vee})_{N_{0}})^{\ord}
 \\
\H^{\textrm{\'et}}_{d}(\baar{X}_{*,U_{*}^{p}U_{*,r}}, \cW)^{\ord}&=\H^{\textrm{\'et}}_{d}(\baar{X}_{*,U_{*}^{p}U_{*,r}}, \cW^{\circ})^{\ord}\ot_{\OO_{L}}L
\eeqq
with respect to the action of $T_{\G_{*}}^{+}$ by $\Up_{t}\ot t$, as defined in \S~\ref{ssec ord}.
The \emph{ordinary completed homology} of $X_{\G_{*}, U^{p}_{*}}$ is 
$$M^{\circ}_{\G_{*}, U_{*}^{p}, W}:=  
\varprojlim_{r } \H^{\textrm{\'et}}_{d}(\baar{X}_{*,U^{p}_{*}U_{*,p,r}}, \cW^{\circ})^{\ord},$$
an $\OO_{L}$-module. It depends on the choice of lattice $W^{\circ}\subset W$, whereas the $L$-vector space 
$$M_{\G_{*}, U_{*}^{p}, W}:=M_{\G_{*}, U_{*}^{p}, W}^{\circ}\otimes_{\OO_{L}}L$$
does not.
When $\cW={\Q_{p}}$ is the trivial local system, we omit it from the notation, thus
$$M_{\G_{*}, U_{*}^{p}}=
M_{\G_{*}, U_{*}^{p}, {\Q_{p}}}.$$






\subsubsection{Independence of weight and Control Theorem} 
For  a $\Z_{p}$-algebra $A$, let  $\cH_{\G_{*}, U^{p}_{*}, p, A}^{\ord}:= A[T_{\G_{*}}]\otimes_{\Z_{p}[T^{+}_{\G_{*},0}]}\Lambda^{\circ}_{\G_{*}, U^{p}_{*}}$. 
For $?= S,\emptyset, {\rm sph}$, consider the $\Lambda_{\G_{*}, U^{p}_{*},A}$-algebra
\begin{align} \label{hecke}
\source{universal-p-adic-gross-zagier-formula}
\cH_{\G_{*}, U^{p}_{*},A}^{?, \ord}:= \cH^{?}_{\G_{*}, U^{p}_{*}, A} \otimes_{A} \cH_{\G_{*}, U^{p}_{*}, p, A}^{\ord}.
\end{align}
For every irreducible algebraic representation $W$ over $L$ and $\OO_{L}$-algebra $A$, the space $M^{\circ}_{\G_{*}, U_{*}^{p}, W}\otimes A$ is  a module over $\cH^{\ord}_{\G_{*}, U^{p}_{*},A}$, where $[t]\in A[T_{\G_{*}}^{+}]$ acts by the double coset operator $\Up_{t}$. 

The base ring $A$ will be omitted from the notation when it can be understood from the context.

Let $U_{*, r}=U^{p}_{*}U_{*, r, p}$ be as in \S~\ref{sec: not} and let $X_{*, r}:=X_{*, U_{*,r}}$. 










\begin{prop}
\source{universal-p-adic-gross-zagier-formula}
\notready \label{free indep} Let $W$ be an irreducible right algebraic representation of $\G_{*/\Q_{p}}$ over $L$, $\cW$ the corresponding local system. Then:
\source{universal-p-adic-gross-zagier-formula}
\begin{enumerate} 
\item If
 $\G_{*}=\G$, $\H$, then   $M_{\G_{*}, U_{*}^{p},W}^{\circ}$ is a projective  $\Lm^{\circ}_{\G_{*}, U^{p}_{*}}\otimes \OO_{L}$-module of finite type. For all  of the groups $\G_{*}$, the  $\Lm^{\circ}_{(\G\times \H)'}\otimes \OO_{L}$-module $M_{(\G\times \H)',K^{p}, W}^{\circ}$  is of finite type, and $M_{(\G\times \H)',K^{p},W}$  is a projective $\Lambda_{(\G\times \H)'}\otimes L$-module of finite type.
\item \label{indep}  We have natural $\cH_{\G_{*}, U^{p}_{*}}^{\ord}$-equivariant isomorphisms  
\source{universal-p-adic-gross-zagier-formula}
\beq\label{iotaW}
\source{universal-p-adic-gross-zagier-formula}
j_{W}\colon M_{\G_{*}, U_{*}^{p}}\otimes \OO_{L}\cong M_{\G_{*}, U_{*}^{p},W}.  \eeq
\item Consider
\beq\label{MGr} 
\source{universal-p-adic-gross-zagier-formula}
M_{\G_{*}, U^{p}_{*}, W,r}:= M_{\G_{*}, U^{p}_{*}}\otimes_{\Lambda_{\G_{*}, U^{p}_{*}}}\Lambda_{\G_{*}, U^{p}_{*}, W, r}.
\eeq
There is a natural $\cH_{\G_{*}, U^{p}_{*}}^{\ord}$-equivariant isomorphism 
$$M_{\G_{*}, U_{*}, W, r} \cong \H_{d}(\baar{X}_{*, r}, \cW)^{\ord}.$$
 \end{enumerate}
\end{prop}
\begin{proof} We first treat part 1 when $W=\Q_{p}$. Then we will deal with part 2, which implies that part 1 holds for any $W$.

If $\G_{*}=\G$, the result is proved  in \cite[Thoerem 1.2, cf. also Remark 1.1]{hida luminy}.
 If $\G_{*}=\H$, $U^{p}_{*}=V^{p}$, then the module under consideration is  isomorphic to   $\Z_{p}
\llb E^{\times}\bks E_{\A^{\infty}}^{\times}/V^{p}\rrb$, which is finite free over $\Lm^{\circ}_{\H,V^{p}}=\Z_{p}\llb \baar{\OO_{E}^{\times }\cap V^{p}}\bks \OO_{E, p}^{\times}  \rrb$ as $ \baar{\OO_{E}^{\times }\cap V^{p}}\bks \OO_{E, p}^{\times}\subset E^{\times}\bks E_{\A^{\infty}}^{\times}/V^{p}$ is a subgroup of finite index.

If $\G_{*}=\G\times \H$, by the K\"{u}nneth formula we have $M^{\circ}_{\G\times \H, U^{p}\times V^{p}}= M^{\circ}_{\G, U^{p}}\hat{\otimes }    M^{\circ}_{\H, V^{p}}$, which by the previous results is a finite type projective module over $\Lm^{\circ}_{\G\times \H} = \Lm^{\circ}_{\G}\hat{\otimes} \Lm^{\circ}_{\H}$. Finally, if $\G_{*}=(\G\times \H)'$ and $K^{p}$ is the image of $U^{p}\times V^{p}$,  by the description of $Z_{K}$ in \eqref{z quot} we have 
\begin{align}\label{mgh}
\source{universal-p-adic-gross-zagier-formula}
M_{(\G\times\H)', K^{p}}^{\circ}= (M^{\circ}_{\G\times \H, U^{p}\times V^{p}})  /  (F_{\A^{\infty}}^{\ts}/ F^{\ts}\cd ((U^{p}\cap F_{\A^{p\infty}}^{\ts}) \cap (V^{p}\cap F_{\A^{p\infty}}^{\ts})))
\end{align}
As $M^{\circ}_{\G\times\H, U^{p}\times V^{p}}$ is a projective $\Lm^{\circ}_{\G\times \H, U^{p}\times V^{p}}$-module of finite type, the quotient
$M^{\circ}_{\G, \times\H, U^{p}\times V^{p}}/{F, p}^{\times} = M_{\G, \times\H, U^{p}N_{\G,{0}}\times V^{p}} \otimes_{\Lm^{\circ}_{\G\times \H}}\Lm^{\circ}_{(\G\times \H)'}$ is a projective  $\Lm^{\circ}_{(\G\times \H)', K^{p}}$-module of finite type, and $M^{\circ}_{(\G\times\H)',K^{p}}$ is its quotient by the free action of the finite group 
$F_{\A^{\infty}}^{\ts}/ F^{\ts}\cd  F_{p}^{\ts}((U^{p}\cap F_{\A^{p\infty}}^{\ts}) \cap (V^{p}\cap F_{\A^{p\infty}}^{\ts}))).$
After inverting $p$, the quotient map admits a section, hence  $M_{(\G\times \H)',K^{p}}$ is projective over $\Lambda_{(\G\times \H)'}$. 

We now turn to part 2. As above it suffices to prove the result  when $\G_{*}=\G, \H$.  Let $\G_{*}=\G$,  and suppose that $W=W_{\G, \uw}^{*}$.
Let $W^{\circ}\subset W$  be the lattice of \eqref{homog poly lattice}, $r\geq 1$. We have a $\Lm_{r}$-linear map
\beq\label{jWr}
\source{universal-p-adic-gross-zagier-formula}
j_{W, r}\colon H_{1}(\baar{X}_{r  }, \Z/p^{r}\Z)^{\ord}
\ot_{\Lm_{r}}   W^{\circ,N_{0}}/p^{r}
\to H_{1}(\baar{X}_{r}, \cW^{\circ}/p^{r}\cW^{\circ})
\eeq
induced by cap product\footnote{I am grateful to David Loeffler and Sarah Zerbes for explaining to me this point of view on the Control Theorem.} via the isomorphism of $\Lm_{r}$-modules $ H^{0}(\baar{X}_{r}, 
\cW^{\circ}/p^{r} \cW^{\circ})\cong W^{\circ,N_{0}}/p^{r}$


The maps \eqref{jWr}  are compatible with variation in $r$, and taking limits  we obtain the map \eqref{iotaW}, which Hida  (\cite[\S~ 8]{Hida88},   \cite[Theorem 2.4]{hida luminy}) proved to be an isomorphism; the asserted equivariance  properties are clear from the construction. 







When $\G_{*}=\H$ the construction is similar but  easier, as each $W$ is $1$-dimensional and each of the analogous maps $j_{W,r}$ is an isomorphism.

Finally, we address part 3. As above we may reduce to the case $W=\Q_{p}$ and $\G_{*}=\H$ or $\G_{*}=\G$. The former is clear, and the latter is, in view of part 2, equivalent to the statement
$$M_{\G, U^{p}, \cW}\otimes_{\Lm_{\G, U^{p}}} \Lm_{\G, U^{p}, r}\cong \H_{d}(\baar{X}_{r},  \cW)^{\circ},$$
which is the control theorem of 
 \cite[Theorem 1.2 (3)]{hida luminy}.
\end{proof}


\subsubsection{Ordinary eigenvarieties}

The space  $M^{\circ}_{\G_{*}, U^{p}_{*}}$ has the structure of an $\cH_{\G_{*}, U^{p}_{*}}^{ \ord}$-module (in particular of $\Lambda^{\circ}_{\G_{*}, U^{p}_{*}}$-module), and for $?=\emptyset, {\rm sph}$ and $A$ a $\Z_{p}$-algebra,   we let
$${\bf T}_{\G_{*}, U^{p}_{*}, A}^{{\rm sph}, \ord}$$
   be the image of  $\cH_{\G_{*}, U^{p}_{*}, A}^{{\rm ?}, \ord}$ in $\End_{A}(M^{\circ}_{\G_{*}, U^{p}_{*}}\otimes A)$,  that  is independent of the particular spherical Hecke algebra chosen when $?={\rm sph}$.  When $A=\Z_{p}$ we omit it from the notation.

We may now define 
  $$\cE^{\ord}_{\G_{*}, U_{*}^{p}} :=\Spec {\bf T}_{\G_{*}, U^{p}_{*}, \Q_{p}}^{{\rm sph}, \ord}.$$
  When $\G_{*}=\H$, we will omit the superscript `${\ord}$'. 
 
Let
   $$\kappa_{\G_{*}}\colon \cE^{\ord}_{\G_{*}, U_{*}^{p}}  \to  \mathfrak{W}_{\G_{*}, U^{p}_{*}}  .$$ 
Referring to Definition \ref{clwt}, the zero-dimensional  (ind)-subscheme of classical points (respectively classical points of weight $W$, for an algebraic representation $W$ of $\G_{*}$, respectively classical points of weight $W$ and level $r$) is 
  $$\cE^{\ord, \cl, (W)}_{\G_{*}, U^{p}_{*}, (r)} := \mathfrak{W}_{\G_{*}, U^{p}_{*},(r)}^{\cl, (W)}     \ts_{ \mathfrak{W}_{\G_{*}, U^{p}_{*}}, \kappa} \cE^{\ord}_{\G_{*}, U_{*}^{p}}. $$
  
  We denote by 
  $$\cM_{\G_{*}, U^{p}_{*}}$$
  the sheaf on $ \cE^{\ord}_{\G_{*}, U_{*}^{p}}$ corresponding to $M_{\G_{*}, U^{p}_{*}}$. 
\begin{enonce*}{Notation} \textup{When $\G_{*}= (\G\times \H)'$, we omit the subscripts, thus e.g. for $K^{p}\subset  (\G\times \H)'(\A^{p\infty})$ 
we write
  $$\cE_{K^{p}}^{\ord}:= \cE_{(\G\times \H)', K^{p}}^{\ord}.$$}
\end{enonce*}



By \eqref{mgh},
  ${\bf T}_{K^{p}}^{{\rm sph}, \ord}$ is a quotient of ${\bf T}_{\G\times \H, U^{p}\times V^{p}}^{{\rm sph}, \ord}$ and correspondingly we have a closed immersion
\begin{align}\label{immers}
\source{universal-p-adic-gross-zagier-formula}
\cE_{K^{p}}^{\ord}\into \cE_{\G\times\H, U^{p}\times V^{p}}^{\ord}.
\end{align}
 
 

  
\begin{prop}
\source{universal-p-adic-gross-zagier-formula}
\notready\label{semi-local} 
\source{universal-p-adic-gross-zagier-formula}
The ring $ {\bf T}_{\G_{*}, U^{p}_{*} }^{{\rm sph}, \ord}$ is finite  flat
 over $\Lambda^{\circ}_{\G_{*}, U^{p}_{*}}$, hence semi-local. 
  The  maximal ideals of $ {\bf T}_{\G_{*}, U^{p}}^{{\rm sph}, \ord} $ are in bijection with $G_{{\bf F}_{p}}$-orbits of  characters $\baar{\lambda}\colon   
 {\bf T}_{\G_{*}, U^{p}_{*}}^{{\rm sph}, \ord}\to \baar{\bf F}_{p}$. 
  \end{prop}
 \begin{proof}
The first statement is easy for the group $\H$ and it is  proved in \cite{hida luminy} for the group $\G$. Together they imply the statement for $\G\times \H$ and hence $(\G\times \H)'$.  
As  $ {\bf T}_{\G_{*}, U^{p}}^{{\rm sph}, \ord} $  is topologically finitely generated over $\Z_{p}$, the residue fields of its maximal ideals are finite extensions of ${\bf F}_{p}$; this implies the second statement.
\end{proof}


\begin{lemm}\label{dense}
\source{universal-p-adic-gross-zagier-formula}
Let $W$ be an irreducible algebraic representation of $\G_{*}$. The set $\cE^{\ord, \cl, W}_{\G*, U_{*}}$
 of classical points of weight $W$
 is Zariski-dense in $\cE^{\ord}_{G_{*}, U_{*}^{p}}$.
\end{lemm}
\begin{proof}
By the previous proposition, the map $\kappa_{\G_{*}}$ is finite hence closed. Then the Zariski-density of
 $\cE^{\ord, \cl, W}_{\G*, U_{*}}=\kappa_{\G_{*}}^{-1} ( \mathfrak{W}_{\G_{*}, U^{p}_{*}}^{W})$ reduces to the Zariski-density of $\mathfrak{W}_{\G_{*}, U^{p}_{*}}^{W} \subset \mathfrak{W}_{\G_{*}, U^{p}_{*}}$, which follows from \eqref{alg density}; cf. aso \cite[Lemma 3.8]{SW99}.
 \end{proof}




























  
  
  



\subsubsection{Abelian case}
The structure  of the eigenvariety for the abelian groups $\H:={\rm Res}_{E/\Q}\G_{m}$ and ${\rm Z}={\rm Res}_{F/\Q}\G_{m}$
 is very simple, and we make it explicit  for the group $\H$: we have
$$M_{\H, V^{p}}:= \hat H_{0}(\baar{Y}_{V^{p}})=  \Z_{p}\llb Y_{V^{p}}( \baar{E})\rrb\otimes\Q_{p},$$
the set $ Y_{V^{p}}(\baar{E})$ is a principal homogeneous space for $\Gamma_{E,V^{p}}:=\H(\Q)\bks \H(\A^{\infty})/V^{p}= E^{\times}\bks E_{\A^{\infty}}^{\times}/V^{p}$, and 
$$\cE_{\H, V^{p}}= \Spec \Z_{p}\llb{\Gamma_{E, V^{p}}}\rrb_{\Q_{p}}.$$
(We omit the superscript `$\ord$' which is meaningless here.)
The  classical points 
 $\cE_{\H, V^{p}}^{\cl}\subset \cE_{\H, V^{p}}(\baar\Q_{p})$ parametrise locally algebraic characters of $\Gamma_{E, V^{p}}$.
Finally, the sheaf $\cM_{\H, V^{p}}$ is a trivial line bundle, with actions  by $G_{E}$ given by the universal character 
\begin{align}\label{chiuniv}
\source{universal-p-adic-gross-zagier-formula}
\chi_{\rm univ}\colon G_{E}\to \Gamma_{E, V^{p}}\to\Z_{p}\llb\Gamma_{E}\rrb^{\times},
\end{align} and by  $\H(\A^{\infty})$  given by the inverse  $\chi_{\H, \rm univ}^{-1}$ of the  corresponding automorphic character. We may formally write 
\begin{align}\label{chiHuniv}
\source{universal-p-adic-gross-zagier-formula}
\cM_{\H, V^{p}}= \chi_{\H, \rm univ}^{-1}\otimes \chi_{\rm univ}
\end{align}
as a tensor product of two trivial sheaves, the first one endowed with the $\H(\A^{})$-action only, and the second one with the Galois action by $ \chi_{\rm univ}$ only.%

\subsubsection{Fibres of the sheaves $\cM$} 
Let 
$$(\lambda^{p},\lambda_{p})\colon {\bf T}_{\G, U^{p}}^{{\rm sph}, \ord}\to \OO(\cE_{\G, U^{p}})$$ be the tautological character, 
 and define
  \beq
\begin{aligned}\label{alphauniv}
\source{universal-p-adic-gross-zagier-formula}
\alpha^{\circ}\colon  F_{p}^{\times}&\to \OO(\cE_{\G, U^{p}})^{\times}\\
x&\mapsto \lambda_{p}(\Up_{x}).
\end{aligned}
\eeq
\begin{prop}
\source{universal-p-adic-gross-zagier-formula}
\notready\label{fibreM2} 
\source{universal-p-adic-gross-zagier-formula}
Let $x\in \cE^{\ord, \cl,W}_{\G, U^{p}}$  be a classical point  of weight $W$ and level $r$. 
  %a classical   $p$-adic automorphic representation $\baar\pi_{x}$ of $\G(\A^{\infty})$ over $\baar{\Q}_{p}(x)$ with finite slope, and let $\theta$ be the associated character of $T$.f
 Then:
\begin{enumerate}
\item  Let $U:=U^{p}U_{p,r}$, and let $\cW$ be the local system on $X$ associated with $W$. 
 We have an isomorphism of $\cH^{\ord}_{\G, U^{p}, \Q_{p}(x)}$-modules
$$(\cM_{\G, U^{p}})_{x}\cong  \H_{1}(X_{U, \baar F}, \cW_{\G, \uw})^{\ord}\otimes_{\cH^{\rm sph}_{\G, U^{p}}, \lambda_{x}} \Q_{p}(x).$$
\item 
There exists  a  unique  automorphic representation ${\pi}_{x}$ of $\G(\A^{})$  over ${\Q}_{p}(x)$ of spherical character $\lambda_{x}^{p}$, weight $W$, and unit character $\alpha_{x}^{\circ}$. It satisfies the property
\beq\label{property}
\source{universal-p-adic-gross-zagier-formula}
\pi_{x}^{ {\ord,U^{p}}}
\cong
\Hom_{\Q_{p}(x)[G_{F, S}]}((\cM_{\G, U^{p}})_{x} ,
 \rho_{x})
\eeq
as left $\cH^{\ord}_{\G, U}\otimes \Q_{p}(x)$-modules.
\end{enumerate}\end{prop}
\begin{proof}
Part 1 follows from Proposition \ref{free indep}.3. 
 %The first part follows from  the control theorem [[ add ref to hida luminy, together with the isomorphism of Hecke algebras]]

For part 2, 
fix an embedding $\Q_{p}(x)\into \baar{\Q}_{p}$. By strong multiplicity-one, a representation $\overline{\pi}$ over $\baar{\Q}_{p}$ with character $\lambda_{x}^{p}$ is unique if it exists. By comparing part 1 with \eqref{carayol iso 2}, we find that $\baar{\pi}$ exists and that for such $\baar{\pi}$  property \eqref{property}  holds after base-change to $\baar{\Q}_{p}$. Let $V_{\baar{\pi}}$ be the Galois representation associated with $\pi$ by Theorem \ref{galois for hilbert}, then by looking at Frobenius traces, we see that  $V_{\baar{\pi}}$  has a model $V_{\pi}$ over $\Q_{p}(x)$. It follows again from \eqref{carayol iso 2} that $\pi :=\lim_{U'}\Hom(H_{1}(\baar{X}_{U}, \cW), V_{\pi}) $ is a model of $\baar{\pi}$ that satisfies \eqref{property}.
\end{proof}

In the rest of the paper, we will use without further comment the notation $\pi_{x}^{}$ for the
 representation of $\G(\A^{})$ defined above, for  $x\in \cE_{\G, U^{p}}^{\ord,\cl }$. 


\begin{coro}\label{fibreMad}
\source{universal-p-adic-gross-zagier-formula}
Let $z\in \cE_{K^{p}}^{\ord, \cl}$ be a classical point, and write $z=(x,y)$ via \eqref{immers} and $L:=\Q_{p}(z)$. Let $\omega_{x}$ be the central character of $\pi_{x}$, let $\chi_{\H,y}$ be the  character of $\H(\A^{\infty})$ obtained by specialising $\chi_{\H, \rm univ}$, and let $\chi_{y}$ be the corresponding locally algebraic character of $G_{E,S}$. Write $L:=\Q_{p}(z)$. Then $\omega_{z}:=\chi_{y|F_{\A^{\infty, \times}}}\omega_{x} = \one $, and 
$$\cM_{K^{p},{z}}\cong
(  (\pi_{x}^{ p,\vee, })^{U^{p}}\otimes_{L}\chi_{\H,y}^{-1,p})  \otimes_{L} (V_{x|_{G_{E,S}}}\otimes_{L[G_{E, S]}} \chi_{y})
$$
as $\cH^{K_{S}}_{S, L}\otimes_{L} L[G_{E,S}]$-modules. Here, $G_{E, S}$ acts trivially on the first two tensor factors, and the natural action of 
 $\cH_{\G\ts\H}^{U_{S}\ts V_{S}}$ 
on the first two factors is extended trivially to the whole tensor product, and  it factors to an action of $\cH^{K_{S}}_{S}$.
\end{coro}

\begin{proof}
Let  $\lambda_{x, F}^{p}$ and $\lambda_{y, F}^{p}$ be the restrictions of the characters $\lambda_{x}$, $\lambda_{y}$ to $\Z[F^{\ts}_{\A^{Sp\infty, \times}}/K_{F}^{Sp}]$, and let $\lambda_{F}$ be the restriction of $\lambda_{F,x}\lambda_{F,y}$ to $\Delta'=F^{\ts}_{\A^{Sp\infty, \times}}/K_{F}^{Sp}$. As this groups acts trivially on $M_{K^{p}}$ by \eqref{mgh}, we have $\lambda_{F}=\one$. On the other hand $\lambda_{F}$ equals the restriction of $\omega_{z}$ to $\Delta'$. We deduce that $\omega_{z}$ factors through $C=F^{\ts}_{\A^{\infty}}/F^{\times}F_{\A^{Sp\infty}}K_{F, S}K_{F,p}$ for some open compact $K_{F,p}\subset F_{p}^{\times}$. By weak approximation, $C=\{1\}$, therefore $\omega_{z}=\one$.

By Proposition \ref{fibreM2}, \eqref{chiHuniv}, and \eqref{mgh}, the asserted result holds provided we quotient the right-hand side by the action of $F_{\A^{\infty}}^{\ts}$,  however  this group acts  by $\omega_{z}$, hence trivially.
\end{proof}


\begin{prop}
\source{universal-p-adic-gross-zagier-formula}
\notready\label{etale} The natural map $\kappa\colon \cE^{\ord}_{K^{p}} \to \mathfrak{W}_{K^{p}}$ is \'etale over a neighbourhood of  the classical  points in $\mathfrak{W}^{\cl}_{K^{p}}$. In particular, the space $\cE^{\ord}_{K^{p}}$ is regular at all $z\in \cE^{\ord, \cl}_{K^{p}}$.
\source{universal-p-adic-gross-zagier-formula}
\end{prop}
\begin{proof} As $\kappa $ is finite  flat by Proposition \ref{semi-local}, it suffices to check that the fibre of $\kappa$ over any $x\in \mathfrak{W}^{\cl}_{K^{p}}(\baar{\Q}_{p})$ is isomorphic to  $\baar{\Q}_{p}^{m}$ for some $m$. By \ref{free indep}, \ref{MGr} and \eqref{carayol iso 2}, this fibre  is the spectrum of the image $A_{x}$ of $\mathscr{H}_{K^{p}}^{{\rm sph}, \ord}$
in   $\bigoplus_{z\in \kappa^{-1}(x)} (\Pi_{z}^{K^{p}, \ord})^{\oplus 2}$, where the $\Pi_{z}$ form   a list  of distinct  irreducible representations of $(\G\times \H)'(\A^{p\infty})$ over $\baar{\Q}_{p}$. By strong multiplicity-one, we have $A_{x}\cong \oplus_{z}\baar{\Q}_{p}$. This proves \'etaleness. As $\mathfrak{W}_{K^{p}}$ is regular, we deuce that so is $\cE_{K^{p}}^{\ord}$ in a neighbourhood of classical points.
\end{proof}



 





 
 
 
 




 




 
 




\subsection{Galois representations in families} 
 %  [[If $M'\supset M$ is a field extension, an $M'$-rational representation $\pi'$ is one of the form $\pi\otimes_{M} M'$ for some $M$-rational representation $\pi$. delete if unnecessary]]







We recall the existence of a  universal family of Galois representations over $\X$.
\subsubsection{Representations associated with irreducible pseudocharacters}\lb{3.2new} Recall that an $n$-dimensional   pseudocharacterof $G$ over  a scheme $\X$ is a function $T\colon G\to \OO(\X)$  that `looks like' the trace of an $n$-dimensional representation of $G$ over $\OO(\X)$, see  \cite{rou} for the precise definition. A pseudocharacter $T$ is said to be (absolutely) irreducible  at a point $x\in \X$ if, for any (equivalently, all) geometric point $\baar{x}$ of $\X$ with image $x$, the pullback $\baar{x}^{*}T$  is not the sum of two  pseudocharacters of dimensions $k$, $n-k$ with $0< k<n$. The \emph{irreducibility locus} of $T$ is the set of points of $\X$ at which $T$ is irreducible; it is open  (\cite[\S~7.2.3]{chenevier}).
 
We start by proving that, if $T$ is irreducible, a  representation with trace $T$ is essentially unique when it exists.


\begin{lemm} Let $\X$ be an integral scheme and let  $\cV_{1}$, $\cV_{2}$ be vector bundles of rank $n>0$ over $\X$. Suppose that there is an isomorphism $F\colon \End_{\OO_{\X}}(\cV_{1})\to \End_{\OO_{\X}}(\cV_{2})$. Then there is an invertible $\OO_{\X}$-module $\calL$ and an isomorphism  
$$g \colon \cV_{1}\cong \cV_{2}\otimes \calL$$ inducing $F$ in the sense that 
$F(T)\otimes{\rm id}_{\calL} =g Tg^{-1}$
for all sections $T$ of $\End_{\OO_{\X}}(\cV_{1})$.
\end{lemm}
\begin{proof}
By \cite[Ch. IV]{knus}, any automorphism of an Azumaya  algebra (such as $ \End_{\OO_{\X}}(\cV_{i})$) is Zariski-locally inner.  Therefore there exists an open cover $\{U_{i}\}$ of $\X$ and isomorphisms $g_{i}\colon \Gamma(U_{i}, \cV_{1})\to \Gamma(U_{i}, \cV_{2})$ such that 
\begin{align}
\label{pippo}
\source{universal-p-adic-gross-zagier-formula}
F(T)=g_{i}Tg_{i}^{-1}\end{align} for all $T \in \End_{\OO_{\X}(U_{i})}(V_{1})$. Let $U_{ij}:=U_{i}\cap U_{j}$ and
\begin{align}\label{cij}
\source{universal-p-adic-gross-zagier-formula}
c_{ij}:=g_{i}^{-1}g_{j},
\end{align}
an automorphism of $\cV_{1}$ over $U_{ij}$.  By \eqref{pippo},  $c_{ij} $ commutes with every $T \in \End_{\OO_{\X}(U_{ij})}(\cV_{1})$, hence it is a scalar in $\OO_{\X}(U_{ij})^{\times}$. One verifies easily that the $c_{ij}$ form a cocycle in $H^{1}(\X,\OO_{\X}^{\times})$. Let $\calL$ denote the associated invertible sheaf, which is trivialised by  the cover $\{U_{i}\}$. Then we may view  $g_{i}\colon \Gamma(U_{i}, \cV_{1})\to \Gamma(U_{i}, \cV_{2}\otimes \calL)$. By \eqref{cij}, the $g_{i}$ glue to the desired isomorphism $g\colon \cV_{1}\cong \cV_{2}\otimes \calL$.
\end{proof}

\begin{lemm}\label{uniq}
\source{universal-p-adic-gross-zagier-formula}
Let $\X$ be an integral scheme and $\mathscr{T}\colon G_{F, S}\to \OO(\X)$ an irreducible pseudocharacter of dimension $n$. Let $\cV_{1}$, $\cV_{2}$ be representations of $G_{F,S} $ with trace $\mathscr{T}$. Then
  there exist    a  line bundle $\mathscr{L}$ with trivial Galois action and a $G_{F, S}$-equivariant isomorphism
 $$\cV_{1}\cong \cV_{2} \otimes \mathscr{L}.$$
\end{lemm}
\begin{proof}
Write $G=G_{F,S}$ and let $\calA:=\OO_{\X}[G]/\Ker(\mathscr{T})$. By \cite[Theorem 5.1]{rou}, $\mathscr{A}$ is an Azumaya algebra of rank $4$. By \cite[Corollary 2.9]{saltman}, 
the two natural injective maps $\alpha_{i}\colon \calA\to {\End}_{\OO_{\X}}(\cV_{i})$ are isomorphisms. Then we conclude by the previous lemma.
\end{proof}

\subsubsection{Galois representations in ordinary families} We prove the analogue in Hida families of Theorem \ref{galois for hilbert}.
\begin{lemm}\label{exist rhobar}
\source{universal-p-adic-gross-zagier-formula}
Let $\overline{\lambda}\colon {\bf T}_{\G, U^{p}}^{\rm sph,ord} \to \baar{\bf  F}_{p}$ be a character. Then there is a unique semisimple representation $\rhobar\colon G_{F,S}\to \GL_{2}(\baar{\bf F}_{p})$
such that $\Tr(\rhobar({\rm Fr}_{v}))=q_{v}^{-1}\baar{\lambda}(T_{v})$  for all $v\notin S$.
\end{lemm}
\begin{proof}
The existence follows by lifting $\baar{\lambda}$ to the character $\lambda_{x}$ associated with a classical point $x$ (that is possible thanks to Lemma \ref{dense}), then taking the semisimplification of the reduction modulo $p$ of a lattice in the representation  $\rho_{x}:=\rho_{\pi_{x}}$ of Theorem  \ref{galois for hilbert};   the uniqueness is a consequence of   the Brauer--Nesbitt theorem.
\end{proof}

By Proposition \ref{semi-local} we may decompose ${\bf T}_{\G, U^{p}}^{{\rm sph}, \ord}\cong
 \prod_{\mathfrak{m}}
 {\bf T}_{\G, U^{p}, \frak m}^{{\rm sph} \ord}$
 and consequently  the generic fibre of the associated schemes also decomposes as 
\begin{align}\label{decomp}
\source{universal-p-adic-gross-zagier-formula}
\cE_{\G, U^{p}}^{\ord} \cong \coprod
\cE_{\G, U^{p}, \frak m}^{\ord}.
\end{align}
We will  say that a connected subset  $\X \subset \cE_{\G, U^{p}}^{\ord}$  has residual representation $\rhobar\colon G_{F}\to \GL_{2}(\baar{\bf F}_{p})$ if $\X$ is contained in some $\cE_{\G, U^{p}, \frak m}^{\ord}$ such that  the character $\lambda_{\frak m}\otimes_{{\bf F}_{p}(\frak m)}\baar{\bf F}_{p} $ associated with $\frak m$ is the character of $\rhobar$.


\begin{prop}
\source{universal-p-adic-gross-zagier-formula}
\notready\label{galois2}
\source{universal-p-adic-gross-zagier-formula}
  Let $\X_{\G}$ be  an irreducible component of $\cE_{\G}$  (that is, a {Hida family}).
 Then there exist:
 \begin{itemize}
 \item
an open subset $\X_{\G}'\subset \X_{\G}$ containing $\X_{\G}^{\rm cl}:=\X_{\G}\cap\cE_{\G, U^{p}}^{\cl}$;
\item a locally free $\OO_{\X_{\G}'}$-module $\cV_{\G}$ of rank two over $\X'_{\G}$, such that $$\cV_{\G, x}\cong V_{\pi_{x}}$$ for all $x\in \X_{\G}^{\cl}$;
\item  a  filtration
\begin{gather}\label{decatv} 0\to \cV_{\G, v}^{+}\to \cV_{\G, v}\to\cV_{\G,v}^{-}\to0,\end{gather}
\source{universal-p-adic-gross-zagier-formula}
where the $\cV_{\G,v}^{\pm}$ are locally free $\OO_{\X'_{\G}}$-modules of rank $1$, and $G_{F_{v}}$ acts on $\cV_{\G, v}^{+}$ by the character associated, via local class field theory, with the character 
\beq\lb{alpha gal}
\alpha_{|F_{v}^{\ts}}^{\circ} \langle \ \rangle_{F_{v}}\eeq
deduced from \eqref{alphauniv}.
 \end{itemize}

The representation $\cV_{\G}$ is uniquely determined up to automorphisms and twisting by line bundles with trivial Galois action.
\end{prop}
The result is due to Hida and Wiles (\cite[\S~ 3.2.3]{fouquet} and references therein), except for the existence of \eqref{decatv} when  the residual Galois representation of $\X_{\G}$ is reducible.
\begin{proof}
Let $\mathscr{T}\colon G_{F, S}\to \OO(\X_{\G})$ be the pseudocharacter defined by $\mathscr{T}({\rm Fr}_{v})=q_{v}^{-1}\lambda(T_{v})$, where $\lambda\colon {\bf T}_{\G, U^{p}}^{\rm sph}\to \OO(\X_{\G})$ is the tautological character. Let $\X_{\G}  ^{\rm irr}\subset \X_{\G}$ be the (open) irreducibility locus. By Theorem \ref{galois for hilbert}, $\X_{\G}^{\cl}\subset \X_{\G}^{\rm irr}$. By Lemma \ref{uniq}, a representation $\cV_{\G}$ is unique up to Galois-trivial twists if it exists. We show existence. 


By \cite[Theorem 5.1]{rou}, $\mathscr{A}:=\OO_{\X_{\G}^{\rm irr}}[G_{F, S}]/ \Ker(\mathscr{T})$ is an Azumaya algebra of rank $4$ over $\cE_{\G'}$ and the natural map 
$$\rho\colon G_{F,S}\to \mathscr{A}^{\times}$$ satisfies ${\rm Tr}\circ \rho= \mathscr{T}$ (where $\Tr$ is the reduced trace of $\mathscr A$). Let $c\in G_{F, S}$ be a complex conjugation; we have an isomorphism $\mathscr A={\mathscr A} (\rho(c)-1)\oplus {\mathscr A}(\rho(c)+1)=:\cV_{+1}\oplus \cV_{-1}$. Each of the $c$-eigen-summands $\cV_{\pm1}$ is a locally free $\OO_{\X_{\G}^{\rm irr}}$-module (since so is $\mathscr A$), whose rank is  $2$: indeed at any classical geometric point  $x\in \X_{\G}^{\cl}(\bC_{p})$, the specialisation $\rho_{x}$ is odd, hence we can pick an isomorphism ${\mathscr A}_{x}\cong M_{2}(\bC_{p})$ sending $\rho_{x}(c)$ to $\smalltwomat 1{}{}{-1}$ from which it is immediate that $\cV_{\pm 1, x}$ has rank $2$; since classical points are dense, we conclude that $\cV_{\pm 1}$ also has rank $2$. 

 Let $\cV_{\G}$ be either of $\cV_{\rm 1, \pm}$. By \cite[Corollary 2.9 (a)]{saltman}, the natural map 
 $${\mathscr A}\to \End_{\OO_{\X^{\rm irr}_{\G}}}(\cV_{\G})$$ 
 is an isomorphism; we view it as an identification to obtain a representation $\rho'$ with trace $\mathscr{T}$. As an irreducible  $2$-dimensional Galois representation over a field is uniquely determined by its trace, the representation $\cV_{\G, x}$ is isomorphic to $V_{\pi_{x}}$. 


We now show the existence of the filtration up  to further restricting the base. Fix a place $v\vert p$ of $F$, and let  $\det_{v}\colon G_{F_{v}}\to \OO(\X_{\G}^{\rm irr})^{\ts}$ be the character giving the action on $\det \cV_{\G,v}$. Let $\cV_{0,v}^{+}$ be the trivial  sheaf $\OO_{\X^{\rm irr}_{\G}}$ with $G_{F_{v}}$-action by the character \eqref{alpha gal}, $\cV_{0,v}:=\cV_{\G,v}$, $\cV_{0,v}^{-}:=(\cV_{0,v}^{+})^{-1}(\det_{v})$.  Finally, for $?=+, -. \emptyset$, let 
$$\cW_{v}^{?}:= \Hom_{\OO_{\X'_{\G}}}(\cV_{0,v}^{+}, \cV_{0,v}^{?}).$$ Then for all $x\in \X_{\G}^{\rm cl}$, by Theorem \ref{galois for hilbert} we have exact sequences 
\beq \lb{cW ex seq}
0\to \cW_{v,x}^{+}= \Q_{p}(x) \to \cW_{x}\to \cW_{v,x}^{-},
\eeq
which we  wish to extend to a neighbourhood of $\X_{\G}^{\cl}$. 
From a  consideration of weights based on 
Theorem \ref{galois for hilbert}, we see that for all $x\in \X_{\G}^{\cl}$,  $H^{0}({F_{v}}, \cW_{v,x}^{-})= H^{2}(F_{v}, \cW_{v,x}^{-})=0$.
 Then from \eqref{cW ex seq} we deduce 
\beq\lb{h2va}
H^{2}(F_{v}, \cW_{v,x})=0\eeq
for all
$x\in \X_{\G}^{\cl }$,
 and from  the Euler--Poincar\'e formula and \eqref{cW ex seq} we deduce that  \beq\lb{consth1}\dim_{\Q_{p}(x)} H^{1}(F_{v}, \cW_{v,x}^{-})= 1+[F_{v}:\Q_{p}], \qquad H^{1}(F_{v}, \cW_{v,x})= 1+2[F_{v}:\Q_{p}]    \eeq 
for all $x\in \X_{\G}^{\cl }.$


By Proposition \ref{torss}.3 below, \eqref{h2va} and \eqref{consth1} imply that  the natural map 
$$H^{0}(F_{v}, \cW)\ot_{\OO_{\X^{\rm irr}_{\G}}}\Q_{p}(x)\to H^{0}(F_{v}, \cW_{v,x})\cong \Q_{p}(x)$$
is an isomorphism for all $x\in \X_{\G}^{\rm cl}$. Hence the sheaf $\calL:= H^{0}(F_{v}, \cW_{v})$ is locally free of rank one  in a neighborhood $\X_{\G}'\subset \X_{\G}^{\rm irr}$ of $\X_{\G}^{\cl}$. Defining 
$$\cV_{\G, v}^{+}:=\calL\ot_{\OO_{\X'_{\G}}} \cV^{+}_{0,v}, $$
the natural map $\cV_{\G, v}^{+}\to \cV_{\G,v|\X_{\G}'}$ is injective, and its cokernel $\cV_{\G,v}^{-}$ has rank one at each $x\in \X_{\G}^{\cl}$. Up to further restricting $\X_{\G}'$, $\cV_{\G,v}^{-}$ is also locally free of rank one. It follows immediately from the construction that the exact sequence 
$$0\to \cV_{\G, v}^{+}\to \cV_{\G, v}\to\cV_{\G,v}^{-}\to 0$$ 
has the asserted properties.
\end{proof}










\begin{prop}
\source{universal-p-adic-gross-zagier-formula}
\notready\label{schur} In the situation of Proposition \ref{galois2},  
\source{universal-p-adic-gross-zagier-formula}
the natural injective map
$$i\colon \OO_{\X'_{\G}}\to \End_{\OO_{\X'_{\G}}[G_{F, S}]}(\cV_{\G})$$
is an isomorphism over an open subset $\X_{\G}''\subset \X_{\G}$ containing $\X^{\cl}$.
\end{prop}
\begin{proof}
By Theorem \ref{galois for hilbert}, $\rho_{x}$ is absolutely irreducible for all $x\in \X_{\G}^{\cl}$. 
 %(since the weight of $x$ is regular. )
We deduce that for each $x\in \X_{\G}^{\cl}$, the map 
$i_{x}$ is an isomorphism. Then we may take for $\X''_{\G}$ the open complement of the support of  ${\rm Coker}(i)$. 
\end{proof}





















  
    















\subsection{Universal ordinary representation and local-global compatibility}\label{3.2}
\source{universal-p-adic-gross-zagier-formula}
The idealised description of what is achieved in this subsection would be to  define a universal ordinary automorphic representation of $\G(\A^{\infty})$  over an irreducible component $\X$ of $\cE_{\G}^{\ord}$; then show that it decomposes as the product of the representations of the local groups $\B_{v}^{\times}$,
for $v\nmid p$,\footnote{The action of $\B_{p}^{\times}$ has already been traded for an action of the torus, subsumed into the $\cE_{\G}$-module structure}     associated to $\cV|_{G_{F,v}}$ by a local Langlands correspondence in  families. 
The definition should be an elaboration of
\begin{align}\label{fakejpi}
\source{universal-p-adic-gross-zagier-formula}
``\Pi_{\G}:=\lim_{U^{p}{}'}\underline{\Hom}_{\OO_{\X}[G_{F, E}]}(\cM_{U^{p}{}'}, \cV_{\G})\text{''}.
\end{align}
For technical reasons, a few modifications are necessary: 
\begin{itemize}
\item the local  Langlands correspondence in families  is not defined for the unit groups of division algebras;\footnote{There is an essential reason for this, namely the possible presence of Schur indices in representations of those groups.}
therefore we  ``remove'' the components at the ramification primes $\Sigma$ of $\B$, in the following way: we consider a component of $\cE$ rather than $\cE_{\G}$, and we take ${H_{\Sg}'}$-coinvariants in an analogue $\Pi$ of \eqref{fakejpi}. For sufficiently large levels, this isolates a local factor of $\Pi$ that is generically free of rank one along  locally distinguished Hida families;
\item in the limit in \eqref{fakejpi}, we fix an arbitrarily large finite set of primes $\Sg'$, disjoint from $\Sigma$ and from $S_{p}$, and we let only the $\Sg'$-component of  $U^{p}{}'$  shrink, so as to get a representation of $\B_{\Sg'}^{\times}$;  
\item 
 we replace the abstractly constructed $\cV=\cV_{\G}\otimes \cV_{\H}$ (where $\cV_{\H}=\chi_{\rm univ}$) by a more geometric incarnation using the sheaf $\cM$ in `new' level (with respect to the chosen irreducible component).
\end{itemize}


We use the correspondence studied in \cite{LLC}, with the caveat that  strictly speaking  the normalisation chosen there differs by the one fixed here in Theorem \ref{galois for hilbert}.1 by a Tate twist. This is only a matter of book-keeping, and in order to avoid excessive notational burden, we do not signal such Tate twists when referring to the results of \cite{LLC} in the rest of this paper.





\subsubsection{Irreducible components}
Let $\X_{\G}\subset \cE^{\ord}_{\G, U^{p}}$ be an irreducible component.  Fix a place $v$ of $F$ not in $\Sigma\cup S_{p}$.
 Recall that the $v$-level of a representation $\pi_{v}$ of $\GL_{2}(F_{v})$ is the smallest $m$ such that $\pi_{v}^{U_{1}(\vpi_{v}^{m})}\neq 0$, where $U_{1}(\vpi_{v}^{m})= \{\smalltwomat abcd \in \GL_{2}(\OO_{F,v})\, :\, c\equiv d-1\equiv 0 \pmod{\vpi_{v}^{m}\OO_{F,v}}\}$.  Let $m_{x,v}$ be the $v$-level of $x\in\X^{\cl}$.
\begin{lemm} The function $x\mapsto m_{x,v}$ is constant on $\X_{\G}^{\cl}$.
\end{lemm}
\begin{proof} 
By \cite{carayol-h}, $m_{x,v}$ equals the conductor of the $G_{F_{v}}$-representation $\cV_{x}$; as all those Galois representations are pure, we may conclude by  \cite[Theorem 3.4]{saha-cond}. 
\end{proof}
 
 We may then define the $v$-level $m_{v}$ of $\X_{\G}$ to be the common value of the $m_{x,v}$ for $x\in \X^{\cl}$.  By the following lemma, it is not restrictive to make the following assumption: \emph{for all $v\notin \Sigma \cup S_{p}$, we have $U_{v}=U_{1}{(\vpi}_{v}^{m_{v}(\X_{\G})})$}. (We say that $\X_{\G}$ is a $v$-new component of $\cE_{\G, U^{p}}^{\ord}$.)
 
\begin{lemm}\label{change level} Let ${}'\X_{\G}\subset \cE^{\ord}_{\G, U^{p}{}'}$ be an irreducible component, and suppose that $U^{p}{}'=\prod_{v\nmid p} U_{v}'$. Let $m_{v}$ be the level of $\X_{\G}$ and let $U^{p}=\prod_{v\nmid p} U_{v}$, with $U_{v}=U_{1}(\vpi_{v}^{m_{v}})\supset U_{v}'$ for all $v\notin \Sigma\cup S_{p}$, and $U_{v}=U_{v}'$.  There exists a unique irreducible component $\X_{\G}\subset \cE^{\ord}_{\G, U^{p}}$ whose image under the natural embedding $\cE^{\ord}_{\G, U^{p}}\subset \cE^{\ord}_{\G, U^{p}{}'}$ is ${}'\X_{\G}$.
\source{universal-p-adic-gross-zagier-formula}
\end{lemm}
\begin{proof} Let $x'\in {}'\X_{\G}^{\cl}$ be any classical point. By \cite{carayol},  its level (that is, the level of $\pi_{x,v}$) is $m_{v}$ if and only if  $\pi_{x,v}$ already occurs in the cohomology of $X_{\baar{F}}$ at $v$-level $m_{v}$, equivalently if and only if (the system of Hecke- and $\Up_{v}$-eigenvalues associated with) $\pi_{x,v}$ occurs in a quotient of $\cM_{U^{p}}$; that is, if $x'$ comes from a point $x$ of $\cE_{\G, U^{p}}$. Let $\X_{\G}\subset \cE^{\ord}_{\G, U^{p}}$ be the irreducible component containing $x$, which is unique by Proposition \ref{etale}. As $\cE^{\ord}_{\G, U^{p}}\subset \cE^{\ord}_{\G, U^{p}{}'}$ are equidimensional of the same dimension, the image of $\X_{\G}$  in $\cE_{\G, U^{p}{}'}$ is an irreducible component, necessarily ${}'\X_{\G}$. 
\end{proof}


We now deal with the level at $\Sigma$.
\begin{lemm}\label{levSig}
\source{universal-p-adic-gross-zagier-formula}
 Let $v\in \Sigma$. There exists a compact open $U_{v}'\subset U_{v}$ such that for every classical point $x'\in \X_{\G}$, we have $\pi_{x',v}^{U'_{v}}=\pi_{x',v}$,
where $\pi_{x,v}$ is the local component at $v$ of $\pi_{x}\otimes \baar{\Q}_{p}$.
\end{lemm}
\begin{proof} 
Fix a classical point $x\in \X_{\G} \subset \cE^{\ord}_{\G, U^{p}}$, and let $U_{v}'\subset U_{v}$ be such that
 $\pi_{x,v}^{U'_{v}}=\pi_{x,v}$. (This will hold for sufficiently small $U_{v}'$ as $\pi_{x,v}$ is finite-dimensional.) We show that $U_{v}'$ satisfies the desired property at all classical $x'\in \X_{\G}$.
Let $\mathfrak{X}_{v/\Q_{p}}$ be the   Bernstein variety of $\GL_{2}(F_{v})$, a scheme over $\Q_{p}$  (see \cite{LLC}, to which we refer for more background). 
By \cite[Theorem 3.2.1]{LLC}, the representation $\cV_{\G}$ of $G_{F_{v}}$ gives a map $f\colon \X_{\G}\to \mathfrak{X}_{v/\Q_{p}}$,
 compatibly with the local Langlands correspondence in the sense that for all $x\in  \X_{\G}$,   $f(x)$ is the point corresponding to the supercuspidal support of the representation $\pi_{x,v}'$ of $\GL_{2}(F_{v})$ over $\Q_{p}(x)$ attached to the representation $\cV_{\G, x}$. Note that for classical points $x$, $\pi_{x,v}={\rm JL}_{v}(\pi'_{x,v} \otimes_{\Q_{p}(x)}\baar{\Q_{p}})$, where  ${\rm JL}_{v}$ is the Jacquet--Langlands correspondence.

 After base-change to $\baar{\Q}_{p}$, we may consider the  finitely many maps $f_{i}\colon \X_{i}\to \mathfrak{X}_{v/\baar{\Q}_{p}}$, where the $\X_{i}$ are the connected components of $\X_{\G,\baar{\Q}_{p}}$. The image of  $f_{i}$ is contained in a connected component $\mathfrak{X}_{i}$ of  $\mathfrak{X}_{v/\baar{\Q}_{p}}$. These components are in bijection with inertial classes of supercuspidal supports for $\GL_{2}(F_{v})$, and  for the class $\sigma=\sigma_i $ of $\frak X_i$ there are three possibilities:
\begin{itemize}
\item $\sigma$ corresponds to  the class of a  supercuspidal representation  $\sigma_0$ of $\GL_{2}(F_{v})$ over $\baar{\Q}_{p}$. In this case, there is an unramified character $\omega\colon F_{v}^{\times}\to \OO(\frak{X}_i)^{\times}$ such that every $y'\in \mathfrak{X}_{i}$ corresponds to $\sigma_0\otimes\omega_{y'}$. Hence for  every classical  $x'\in  \X_{i}$, we have $$\pi_{x',v}\cong {\rm JL}_{v}( \sigma_0\otimes\omega_{f_{i}(x')})={\rm JL}_{v}( \sigma_0)\otimes\omega_{f_{i}(x')}\cong \pi_{x,v}\otimes \omega_{f_{i}(x')}.$$
 As $\omega_{f(x')}$ is unramified, it follows that  $\pi_{x',v}^{U_{v}}=\pi_{x',v}$.
\item $\sigma$ corresponds to the class of the supercuspidal support of ${\rm St}\otimes\omega_{0}$, where ${\rm St}$ is the Steinberg representation and $\omega_{0}\colon F_{v}^{\times}\to \baar{\Q}_{p}^{\times}$ is a character. Then there exist a closed subset $\mathfrak{X}_{i}'\subset \mathfrak{X}_{i}$ and an unramified character $\omega\colon F_{v}^{\times}\to \OO(\mathfrak{X}_{i})^{\times}$ such that  every $y'\in \mathfrak{X}_{i}'$ corresponds to the supercuspidal support of ${\rm St} \otimes\omega_{0}\omega_{y'}$, and such that every $y'\in \mathfrak{X}_{i}-\mathfrak{X}_{i}'$ corresponds to the support of an irreducible principal series representation. It follows that for  every classical $x'\in \X_{i}$,  the image $f_{i}(x')\in \mathfrak{X}_{i}'$ (since $\pi_{x,v}'\otimes \baar{\Q}_{p}$ is in the domain of ${\rm JL}_{v}$),  and that $\pi_{x',v} ={\rm JL}_{v}({\rm St}\otimes\omega_{0}\omega_{f_{i}(x')})=\omega_{0}\omega_{f_{i}(x')}\circ {\rm Nm}$. We conclude as above.

\item no element of the inertial class $\sigma$ is the supercuspidal support of a special or supercuspidal representation. This case is excluded as only those representations are in the image of the Jacquet--Langlands correspondence.
\end{itemize}
\end{proof}


\subsubsection{Galois representation from geometry}
Let $U^{p}\subset \G(\A^{p\infty})$, $V^{p}\subset \H(\A^{p\infty})$ be compact open subgroups. We will consider various  compact open subgroups $U^{p}\subset U^{p}_{*} \subset \G(\A^{p\infty})$, and will correspondingly denote by   $K^{p}_{*}$ be the image of $U^{p}_{*}\times V^{p}$ in $(\G\times \H)'(\A^{p\infty})$. Let $\X$ be an irreducible component of $\cE_{K_{}^{p}}^{\ord}\subset \cE_{\G, U^{p}}^{\ord}\times \cE_{\H, V^{p}}$, and let $\X_{\G}\subset\cE_{\G, U^{p}}^{\ord}$ be the irreducible component such that $\X\subset \X_{\G\ts \H}:=\X_{\G}\ts  \cE_{\H, V^{p}}$. 


Suppose from now on that $\X$ is locally distinguished by $\H'$ (Definition \ref{hida loc dist}).
Let  $\cV_{\G}$ be the $\OO_{\X_{\G}''}[G_{F, S}]$-module   constructed in Proposition \ref{galois2} and Proposition \ref{schur}, and let $\cV_{\H}$ be the universal character $\chi_{\rm univ}$ of $G_{E, S}$   from \eqref{chiuniv}. 
Let 
$$\X^{(0)}:= \X\cap (\X_{\G}'' \ts  \cE_{\H, V^{p}})$$ an open subset,  and consider the $G_{F, E, S}$-representation 
$$\cV':=(\cV_{\G}\boxtimes \cV_{\H})_{|\X^{(0)}}$$

We  define another  sheaf $\cV$  with $G_{F, E, S}$-action,  that will provide a more  convenient and concrete substitute for $\cV'$ on (an open subset of) $\X^{(0)}$.

 Let 
$$U_{0}^{p}{}'=U^{\Sigma p}\prod_{v\in \Sigma } U_{v}',$$ with $U_{v}'$ as in Lemma \ref{levSig}.  Let $K_{0}^{p}{}'=(U_{0}^{p}{}'\times V^{p})F_{\A^{p\infty}}^{\times}/F_{\A^{p\infty}}^{\times}$, and let
$$\cV:=\cM_{K_{0}^{p}{}'}^{{H_{\Sg}'}},$$ viewed a sheaf over 
$\X$.


\begin{lemm}\label{dirsum}
\source{universal-p-adic-gross-zagier-formula}
 The sheaf $\cV$ is a direct summand of $\cM_{K_{0}^{p}{}'}$.
 \end{lemm}
 \begin{proof} The  group ${H_{\Sg}'}=\prod_{v\in \Sigma} E_{v}^{\times}/F_{v}^{\ts}$ acts on the locally free sheaf  $\cM_{K^{p}_{0}{}'}$
through a quotient by an open subgroup. Since $H_{\Sg}'$ is compact, such a quotient is  finite; 
 therefore the inclusion $\cV\subset \cM_{K_{0}^{p}{}'}$ splits. 
 \end{proof}

\begin{prop}
\source{universal-p-adic-gross-zagier-formula}
\notready\label{Xad*} 
\source{universal-p-adic-gross-zagier-formula}
  There is an open subset $\X^{(1)}\subset \X$ containing all classical points  such that $\cV$ is locally free of rank~$2$ 
  along $\X^{(1)}$.  For every  $z=(x,y)\in \X^{ \cl}$  we have
$$\cV_{z}\cong V_{{x}}\otimes \chi_{{y}}$$
as a $G_{F, E, S}$-representation.  
\end{prop}
\begin{proof} 
 By 
Corollary \ref{fibreMad},
for $z=(x,y)\in \X^{\cl}$ we have
$$\cM_{K_{0}^{p},(x,y)}\cong(\pi_{x}^{\vee, p, U_{0}^{ p}}\otimes \chi_{y}^{-1, p})\otimes (V_{x}\otimes \chi_{y})$$
(where the first pair of factors is a representations of $\G\times \H(\A^{p\infty})$ and the second one is a a representation of $G_{E}$).
By Lemma \ref{dirsum}, taking ${H_{\Sg}'}$-invariants commutes with specialisation, and we find that 
\begin{align}\label{old}
\source{universal-p-adic-gross-zagier-formula}
\cV_{z}\cong  (\pi_{x}^{\vee, \Sigma p, U_{0}^{ \Sigma p}}\otimes\chi_{y}^{-1,  \Sigma p})\otimes (\pi_{x, \Sigma}^{\vee}\otimes \chi_{y,\Sigma}^{-1})^{{H_{\Sg}'}} \otimes( V_{x}\otimes \chi_{y}).\end{align}
The first factor is $1$-dimensional by the theory of local newforms, and the second factor is $1$-dimensional by assumption $(\eps_{v})'$.

Since the fibre-rank of $\cV$ is $2$ in the dense set $\X^{\cl}$, there is an open neighbourhood  of this set over which $\cV$ is locally free of rank $2$. 
\end{proof}


\begin{coro} There exist:
 an open subset  $\X^{(2)}\subset \X^{(0)}\cap \X^{(1)}$ containing $\X^{\cl}$ such that 
\beqq \End_{\OO_{\X^{(2)}}[G_{F, E,S}]}(\cV)=\OO_{\X^{(2)}},
\eeqq
 an invertible sheaf $\calL$ over $\X^{(2)}$ with trivial Galois action, and a $G_{F, E, S}$-equivariant isomorphism of sheaves on $\X^{(2)}$
$$\cV\cong \calL\otimes \cV'.$$
\end{coro}
\begin{proof} By Proposition \ref{Xad*} and the construction of $\cV'$, the representations $\cV$, $\cV'$ have a common  trace $\mathscr{T}\colon G_{F, E, S}\to \OO(\X^{(1)}) $. Since this is an irreducible pseudocharacter, the assertions follow
 from Lemma \ref{uniq} and (the argument of) Proposition \ref{schur}.
\end{proof}

\subsubsection{The universal ordinary representation} In what follows, all sheaves $\cM_{K^{p}_{*}}$ will be considered as sheaves over $\X$ (or open subsets of $\X$). Note that, as the action of $H'_{\Sg}$ on $\cM_{\G}$, $ \cM_{\H}$ commutes with the Galois action, the sheaves $\cM_{K^{p}_{*}}$ retain an action of $G_{F, E, S}$.  

We will use the following well-known fact. 
\begin{lemm}\lb{detT} Let $R$ be a ring and let $T\colon M\to N$ be a map of free $R$-modules of the same rank. The set of those $x\in \Spec R$ such that $T\otimes R/\frak{p}_{x}$ is an isomorphism is open in $\Spec R$.
\end{lemm}
\begin{proof} The locus is the complement of $V(\det T)$. 
\end{proof}

In what follows, similarly to \S~\ref{ssec cong}, if `$?$' is any decoration,  $U_{?}^{p}$ is a subgroup of $\G(\A^{p\infty})$, and $V^{p}\subset \H(\A^{p\infty})$ is a fixed subgroup,  we denote by $K_{?}^{p}\subset (\G\ts\H)'(\A^{p\infty})$ the image of $U_{?}^{p}\ts V^{p}$. 
\begin{prop}
\source{universal-p-adic-gross-zagier-formula}
\notready\label{cofinal} Fix a finite set of primes $\Sg'$ disjoint from $\Sigma \cup S_{p}$, such that $U_{v}$ is maximal for all $v\notin {\Sg'}\cup \Sigma\cup S_{p}$, and consider the set $\mathscr{U}$ of subgroups   $U^{p}{}'=\prod_{v\nmid p}U_{v}'\subset U^{p}$  with $U'_v$  as in Lemma \ref{levSig} for all $v\in \Sigma$, and $U_{v}'=U_{v}$ for all $v\notin {\Sg'}\cup \Sigma \cup S_{p}$. (In particular $U_{0}'\in \mathscr{U}$.)
\source{universal-p-adic-gross-zagier-formula}
 \begin{enumerate}
\item
There exists a cofinal   sequence $(U^{p}_{i}{}')_{i\geq 0}\subset \mathscr{U} $, and open subsets $\X_{i}\subset \X^{(2)}\subset \X$ containing $\X^{\cl}$ such that $\X_{j}\subset \X_{i}$ for $i\leq j$, satisfying the following: there are  
 integers $r_{i}$ and $G_{F ,E}$-equivariant  maps
$$T_{i}\colon \cV^{\oplus r_{i}}= (\cM_{ K^{p}_{0}{}'}^{{H_{\Sg}'}} )^{\oplus r_{i}}   \to 
\cM_{K^{p}_{i}{}'}^{{H_{\Sg}'}}$$
that 
 are isomorphisms over $\X_{i}$.
\item
For each 
 $U^{p}{}'\in \mathscr{U}$, there is an open subset $\X_{U^{p}{}'}\subset \X^{(2)}$ containing $\X^{\cl}$ such that 
   the restriction to $\X_{U^{p}{}'}$ of 
\beq\label{piku}
\source{universal-p-adic-gross-zagier-formula}
\Pi^{K^{p}{}', \ord}_{{H_{\Sg}'}}:=\underline{\Hom}_{\OO_{\X}[G_{F, E, S}]}(\cM_{K^{p}{}'}^{{H_{\Sg}'}},\cV) \eeq
is a locally free
$\OO_{\X_{U^{p}{}'}}$-module, and we  have  an isomorphism of locally free sheaves with Hecke- and Galois- actions
 $$\cM_{K^{p}{}'}^{{H_{\Sg}'}} \cong 
( \Pi^{K^{p}{}', \ord}_{{H_{\Sg}'}} )^{\vee} \otimes \cV.$$
Moreover $\Pi^{K^{p}{}', \ord}_{{H_{\Sg}'}}\subset \Pi^{K^{p}{}'', \ord}_{{H_{\Sg}'}} $  for $U^{p}{}''\subset U^{p}{}'$
 via the natural projections  $\cM_{ K^{p}{}''}^{{H_{\Sg}'}}\to\cM_{ K^{p}{}'}^{{H_{\Sg}'}} $.
\item  The $ \cH_{\G\times \H, {\Sg'}}^{K^{p}{}'}$-module  $\Pi^{K^{p}{}', \ord}_{{H_{\Sg}'}}$ is generated by  $\Pi^{K_{0}^{p}{}', \ord}_{{H_{\Sg}'}} $
over $\X_{U^{p}{}'}$.
\item\label{Piord}  For each $z=(x,y)\in\X^{\cl}$, we have  
\source{universal-p-adic-gross-zagier-formula}
$$(\Pi^{K^{p}{}', \ord}_{{H_{\Sg}'}})_{z}\cong (\pi_{x}^{U^{p}{}', \ord})_{{H_{\Sg}'}}\otimes \chi_{y}, $$
with the notation of \eqref{piord}.
\end{enumerate}
\end{prop}
\begin{proof} 
It suffices to prove part 1 for a sequence of subgroups $U_{i}^{p}{}'=\prod_{v\nmid p}U_{i,v}'$ that are $\B_{S}^{\times}$-\emph{conjugate} to a cofinal sequence (if $U_{i}^{p}{}''=g_{i}U_{i}^{p}{}'g_{i}^{-1}$ is cofinal and $(U_{i}^{p}{}', T_{i})$ satisfies the desired condition, then $(U_{i}^{p}{}'', g_{i}^{-1}\circ T_{i} )$ also  satisfies the desired condition).
We thus take any sequence with   $U_{i,v}'=U_{1}(\vpi^{m_{i,v}})$ for $v\notin  {\Sg'}\cup\Sigma \cup S_{p}$, with $m_{i,v} \geq m_{v}$ and such that $\min_{v\in {\Sg'}}m_{i,v}\to \infty$.

Let $r_{i}=\prod_{v}(1+m_{i,v}-m_{v})$. By the local theory of oldforms of   \cite{casselman} and the isomorphisms \eqref{old} and
\begin{align}\label{old2}
\source{universal-p-adic-gross-zagier-formula}
 (\cM_{ K^{p}_{i}{}'})^{ {H_{\Sg}'}}_{z}\cong  (\pi_{x}^{\vee, \Sigma p, U_{i}^{ \Sigma p}{}'}\otimes\chi_{y}^{-1,  \Sigma p})\otimes (\pi_{x, \Sigma}^{\vee}\otimes \chi_{y,\Sigma}^{-1})^{{H_{\Sg}'}} \otimes( V_{x}\otimes \chi_{y}),
\end{align}
 there are Hecke operators $$T_{v, j_{v}} \colon \cM_{ K^{p}_{0}{}'}^{{H_{\Sg}'}}     \to 
\cM_{K^{p}_{i}{}'}^{{H_{\Sg}'}}$$ such that the map $T_{i}:=\prod_{v\in S} \oplus_{j_{v}} T_{i, v, j_{v}} $ is an isomorphism after specialisation at any $z$ in the dense set $ \X^{\rm  cl}$. Hence $T_{i}$ is an isomorphism in an open neighbourhood $\X_{i}$ of $ \X^{\cl}$ (which we possibly shrink to make sure it is contained in $\X^{(2)}$).   Together with Lemma \ref{detT}, this concludes the proof of part 1. Part 2 is a consequence of part 1 and the absolute irreducibility of $\cV$, in the special case  $U^{p}{}'=U^{p}_{i}{}'$, with $\X_{U^{p}{}'}=\X_{i}$. The general case is deduced from the special case: if $U^{p}{}'\subset U_{i}^{p}{}'$, let $\X_{U^{p}{}'}:=\X_{i}$ and take  on both sides the locally free summands  consisting of $U^{p}{}'$-invariants (for the first assertion) or coinvariants (for the second assertion). For part 3,  we may again reduce to  the special case $U^{p}{}'=U^{p}_{i}{}'$; then the space $\Pi^{U_{i}^{p}{}'}$  is generated by the images of the transposes of various ``oldforms'' degeneracy maps $T_{i}$ from part 1, that are elements of the Hecke algebra $\cH_{\G, S}^{U^{p}{}'}$. Finally, part \ref{Piord} follows from \eqref{old2}. 
\end{proof}




\begin{defi}[Universal ordinary representation]
\source{universal-p-adic-gross-zagier-formula}
\notready 
Let $\mathscr{U}$ be as in Proposition \ref{cofinal}, and fix an arbitrary 
$ U' \in \mathscr{U}$. Let 
$$\X^{(3)}:=\X_{U^{p}{}'}$$
be as in Proposition \ref{cofinal}, and 
let 
$$\Pi^{K^{p}{}', \ord}_{{H_{\Sg}'}}:=\underline{\Hom}_{\OO_{\X^{(3)}}[G_{F, E, S}]}(\cM_{K^{p}{}'}^{{H_{\Sg}'}},\cV) $$ as in \eqref{piku}.
The \emph{universal ordinary representation}
 $$\Pi^{K^{Sp}, \ord}_{{{H_{\Sg}'}}}\subset
 {}' \Pi^{K^{Sp}, \ord}_{{H_{\Sg}'}}:=\varinjlim_{U^{p}{}''\in \mathscr{U}} \Pi^{K^{p}{}'', \ord},$$
 is the $\OO_{\X^{(3)}}[(\B_{{\Sg'}}^{\times}\times E_{{\Sg'}}^{\times})/F_{{\Sg'}}^{\ts}]$-submodule
  generated by $\Pi_{{H_{\Sg}'}}^{K^{p}{}', \ord}$.  
  \end{defi}








\subsubsection{Local-global compatibility}
The next theorem describes $\Pi^{K^{{S}p}{}', \ord}_{{{H_{\Sg}'}}}$, as a sheaf with an action by $\B_{{\Sg'}}^{\times}\times E^{\Sigma\infty,\times}$,  in terms of the local Langlands correspondence in  families of \cite{LLC}, denoted  by 
$$\cV_{\G}\mapsto \pi_{\G, {\Sg'}}(\cV_{\G}) .$$
This correspondence attaches, to any  family $\cV_{\G}$ of representations of $\prod_{v\in {\Sg'} } G_{F_{v}}$ on  a rank-$2$ locally free sheaf over a Noetherian  scheme   $\Y/\Q$, a   family of representations of $\GL_{2}(F_{{\Sg'}})$ on a torsion-free sheaf over $\Y$.  The latter  representation is \emph{co-Whittaker} in the sense of \cite[Definition 4.2.2]{LLC}; in particular it admits a unique Whittaker model. 



\begin{theo}[Local-global compatibility]\label{LGC} 
\source{universal-p-adic-gross-zagier-formula}
Let 
$$\pi_{\G,{\Sg'}}(\cV_{\G})$$
be the representation of $\GL_{2}(F_{\Sg'})$ over $\X^{(3)}$ associated with $\cV_{\G}$ by the local Langlands correspondence in  families for $\GL_{2}(F_{{\Sg'}})$ of \cite{LLC};  let $\chi_{\H, {\rm univ}, {\Sg'}}$ be
 the pullback to $\X^{(3)}$ of the sheaf  $\chi_{\H, \rm univ}$  of \eqref{chiHuniv}, with the $\H(\A^{\infty})$-action restricted to $E_{{\Sg'}}^{\times}$.

Then there  exist an open subset $\X^{(4)}\subset \X^{ (3)}\subset \X$ containing $\X^{\cl}$,  a line bundle $ \Pi^{K^{Sp}{}',S,  \ord}_{{H_{\Sg}'}}$ over $\X^{ (4)}$, and an isomorphism of $\OO_{\X^{(4)}}[\GL_{2}(F_{{\Sg'}})\times E_{{\Sg'}}^{\times}]$-modules
$$
\Pi^{K^{{S}p}{}',  \ord}_{{{H_{\Sg}'}}}\cong 
(
\pi_{\G,{\Sg'}}(\cV_{\G})\otimes_{\OO_{\X^{(3)}}} \chi_{\H, {\rm univ}, {\Sg'}}
)
  \otimes 
  \Pi^{K^{{S}p}{}',{S},  \ord}_{{H_{\Sg}'}}. $$
\end{theo}
\begin{proof}
For $*=\emptyset, '$,  consider
$${}^{*}\pi_{\G, {\Sg'}}':=\underline{\Hom}_{\OO_{\X^{(3)}}[E_{S}^{ \times}]}  (\chi_{\H, {\rm univ},S}, {}^{*}\Pi^{K^{Sp}, \ord}_{{{H_{\Sg}'}}}),$$
a torsion-free sheaf over $\X^{(3)}$ with action by $\B_{{\Sg'}}^{\times}=\GL_{2}(F_{{\Sg'}})$. There is an obvious isomorphism
\begin{align}\label{subst}
\source{universal-p-adic-gross-zagier-formula}
{}^{*}\Pi^{K^{{S}p}, \ord}_{{H_{\Sg}'}} \cong {}^{*}\pi_{\G,{\Sg'}}'\otimes
 \chi_{\H, {\rm univ}, {\Sg'}}.
\end{align}
 
 
 By 
 Proposition \ref{cofinal}.\ref{Piord}
 and the local freeness of  each
  ${}'\Pi^{K^{p}{}''}_{{H_{\Sg}'}}$ near $\X^{\cl}$, the fibre of  
  $({}'\pi_{\G, S}')^{U_{S}'}$ at any $z=(x,y)\in \X^{\cl}$ equals   $\pi_{\G,S}(\cV_{\G,x})^{U_{S}''}$;    by Proposition \ref{cofinal}.3 the same is true if one replaces $({}^{}\pi'_{\G, {\Sg'}})^{U_{{\Sg'}}''}$ by the submodule $(\pi'_{\G, {\Sg'}})^{U_{S}''}$.
In conclusion, taking the limit over $U^{p}{}''\in \mathscr{U} $ we find that the smooth, finitely generated, admissible  $\OO_{\X^{(3)}}[\GL_{2}(F_{{\Sg'}})]$-module $\pi'_{\G,{\Sg'}}$ satisfies 
$$\pi'_{\G,{\Sg'}, (x,y)}\cong \pi_{\G, {\Sg'}}(\cV_{\G,x}).
$$
for all $(x, y)\in \X^{\cl}$. 
Then by \cite[Theorem  4.4.3]{LLC},  there exist an open subset $\X^{(4)}\subset \X^{(3)}$ containing $\X^{\cl} $ and  a line bundle that we denote by $ \Pi^{K^{Sp},S, \ord }_{{H_{\Sg}'}}$, such that 
$$\pi'_{\G, {\Sg'}}\cong  \pi_{\G, {\Sg'}}(\cV_{\G})\otimes \Pi^{K^{Sp},S, \ord}_{{H_{\Sg}'}}.$$
Substituting in \eqref{subst} gives the desired result.
\end{proof}
